\documentclass[11pt]{article}
\usepackage{amsmath,amssymb,amsthm}
\usepackage[margin=1.1in]{geometry}
\usepackage{booktabs}
\usepackage[colorlinks=true,linkcolor=blue,citecolor=blue,urlcolor=blue]{hyperref}

\newtheorem{theorem}{Theorem}[section]
\newtheorem{lemma}[theorem]{Lemma}
\newtheorem{proposition}[theorem]{Proposition}
\newtheorem{corollary}[theorem]{Corollary}
\newtheorem{conjecture}[theorem]{Conjecture}
\theoremstyle{definition}
\newtheorem{definition}[theorem]{Definition}
\newtheorem{remark}[theorem]{Remark}

\newtheorem*{thmA}{Theorem A (frontier-width lower bound)}
\newtheorem*{thmAp}{Theorem A$'$ (robustness under disjunction and adaptive stopping)}
\newtheorem*{thmB}{Theorem B (exact growth rate)}
\newtheorem*{thmC}{Theorem C (Sturmian letter factorization)}
\newtheorem*{thmD}{Theorem D (standard-word block recursion)}
\newtheorem*{thmE}{Theorem E (non-regularity of the frontier)}
\newtheorem*{thmF}{Theorem F (no bounded dyadic one-step Lyapunov function)}
\newtheorem*{thmG}{Theorem G (cylinder-local realization)}

\DeclareMathOperator{\ord}{ord}
\newcommand{\N}{\mathbb{N}}
\newcommand{\Z}{\mathbb{Z}}
\newcommand{\R}{\mathbb{R}}
\newcommand{\ceil}[1]{\lceil #1\rceil}
\newcommand{\floor}[1]{\lfloor #1\rfloor}

\title{Certificate complexity of uniform descent\\ for the $3x+1$ problem}

\author{Hikaru Wakaura\thanks{QuantScape Inc., 4-11-18,
Manshon-Shimizudai, Meguro, Tokyo, 153-0064, Japan. Corresponding author;
email: \texttt{hikaruwakaura@gmail.com}. 
% TODO(author): add ORCID iD at submission.
}}
\date{\today}
 
\begin{document}

\maketitle

\begin{abstract}
We formalize the most natural proof strategy for the $3x+1$ (Collatz)
problem --- partition the integers into congruence classes and certify
uniform descent on each class by an affine identity --- as a proof system
of dyadic affine descent certificates, and we prove quantitative lower
bounds on its complexity. In every correct system the certificates deeper
than $k$ must occupy all $c(k)$ climbing classes mod $2^k$, and truncation
at depth $K$ leaves uncovered density $c(K)/2^K$; $c(k)$ grows exactly like
the binary entropy $H(\log_3 2)=0.9500$ per digit, with non-asymptotic
two-sided constants. The frontier is unchanged under disjunctions of
congruence classes of arbitrary moduli and adaptive stopping times
(Terras's coefficient stopping time), and an elementary count bounds the
potential counterexamples up to $N$ by $O(N^{H(\log_3 2)})$. The frontier
is driven by the Sturmian word of slope $\log_3 2$ and is not regular.
The counting theorem, the sandwich, and the core lemmas are
machine-verified in Lean~4.
\end{abstract}
 
\noindent\textbf{Keywords:} $3x+1$ problem; Collatz conjecture; proof
complexity; descent certificate; stopping time; Sturmian word; linear
forms in logarithms; machine-verified proof.
 
\noindent\textbf{2020 Mathematics Subject Classification:} Primary 11B83;
Secondary 68R15, 03F20, 11J86.

\noindent\textbf{Word count:} approximately 6{,}600 words.

\section{Introduction}\label{sec:intro}

The $3x+1$ problem (Collatz conjecture) asserts that iterating
$C(n)=n/2$ for even $n$ and $C(n)=3n+1$ for odd $n$ eventually reaches $1$
from every positive integer. Despite its elementary statement, the problem
has resisted every known technique; see the surveys of Lagarias
\cite{Lagarias1985,Lagarias2010} for its history. Verified convergence now
extends to all $n<2^{68}$ \cite{Barina2021}; almost every integer (in
natural density) descends below itself, by Terras \cite{Terras1976} and
Everett \cite{Everett1977}; and Tao \cite{Tao2022} proved that almost all orbits (in logarithmic
density) attain almost bounded values. Concurrently with and independently
of the present work, Sharpe \cite{Sharpe2026} has machine-verified in
Lean~4 a no-go theorem for \emph{bounded-window} descent certificates
together with Terras's theorem in finite form; the relation to our
results is discussed in Remarks~\ref{rem:thmA}, \ref{rem:countnum} and
\ref{rem:sturmrelated}.

Faced with such a problem, the strategy that nearly every newcomer --- and
many published attempts --- adopt is the same folk recipe: split the
integers into congruence classes mod $2^k$, prove on each class an affine
identity forcing descent below the starting point, and conclude by strong
induction. The recipe succeeds on a set of classes of density approaching
$1$, and this is precisely the content of the classical density theorems.
However, the recipe never seems to terminate: no matter how deep the case
analysis is pushed, a residue of unresolved classes survives, and it has
remained unquantified \emph{whether this is an artifact of particular
implementations or a provable property of the strategy itself}.

Here we prove that it is a provable property of the strategy, and we
determine its exact cost. We formalize the recipe as a proof system of
\emph{dyadic affine descent certificates} (Definition~\ref{def:cert}),
flexible enough to cover finite exceptional sets, refinements by odd moduli,
and arbitrary mixtures of depths, and we establish the following.

\begin{enumerate}
\item[(i)] \emph{A frontier-width lower bound against all systems}
  (Theorem~A): in every correct certificate system, the certificates of
  dyadic depth exceeding $k$ must occupy \emph{every} one of the $c(k)$
  \emph{climbing classes} mod $2^k$, and any family of certificates of
  depth $\le K$ leaves uncovered a set of density $c(K)/2^K$. In
  particular no finite certificate system exists, and the cost of covering
  all but density $\varepsilon$ is depth $\log_2(1/\varepsilon)/(1-H(\theta))$
  up to lower order. The bound is robust (Theorem~A$'$): allowing each
  certificate to certify descent at a \emph{variable} stopping time, over a
  \emph{disjunction} of congruence classes of arbitrary moduli, leaves the
  frontier unchanged --- adaptive stopping collapses to the canonical
  first-descent time (the horizon-bounded form of Terras's coefficient
  stopping time), and each constituent class can serve at most one
  climbing class.
\item[(ii)] \emph{The exact growth rate} (Theorem~B): the non-asymptotic
  sandwich
  $\binom{k}{\ceil{\theta k}}/k \le c(k)\le (k+1)\binom{k}{\ceil{\theta k}}$,
  $\theta=\log_3 2$, whence
  $\gamma=\lim_k \tfrac1k\log_2 c(k) = H(\theta) = 0.9499555272\ldots$
  The lower bound rests on an elementary rotation lemma of independent
  interest.
\item[(iii)] \emph{Fine structure} (Theorems~C--E): the climbing condition
  factorizes over the Sturmian word of slope $\theta$; block transfer
  matrices obey the standard-word recursion governed by the continued
  fraction of $\log_2 3$; and the language of climbing words is not
  regular, so the certificate frontier admits no finite-automaton
  compression.
\item[(iv)] \emph{Machine verification}: Theorems~A and B are verified
  in Lean~4 --- A in an abstract certificate-system formulation (as an
  injection from climbing classes into deep certificates), the two-sided
  binomial sandwich of B against Mathlib --- and so are the core arithmetic lemmas behind
  Theorem~A$'$ (the climbing towers $-1\bmod 2^m$, the affine identity
  with its remainder bound, both halves of the adaptive collapse lemma,
  and the persistence step), all from the standard axioms alone
  (Section~\ref{sec:artifacts}).
\item[(v)] \emph{A verified computational artifact}: a depth-$20$
  certificate system (4404 certificates, exact cover of $\N\setminus\{1\}$
  up to the climbing classes, density $97.39\%$) checked by an independent
  symbolic verifier.
\item[(vi)] \emph{Local obstructions} (Theorems~F--G): no bounded
  syndrome correction $\psi(n\bmod 2^d)$ makes $\log_2 n+\psi$ a one-step
  Lyapunov function, and finite-window constraint hierarchies cannot
  exclude climbing statistics, because the all-ones statistics are realized
  by genuine integer orbit segments at every window length.
\item[(vii)] \emph{A power-saving counterexample count}
  (Corollary~\ref{cor:count}): unconditionally, the number of $n\le N$
  whose orbit never descends below $n$ is $O(N^{H(\theta)})$, by an
  elementary argument: the counting exponent is the certificate growth
  rate $H(\theta)=0.9500$ itself.
\end{enumerate}

We stress an explicit non-claim. The density exponent $1-H(\theta)$
of the unresolved classes is classical: it is implicit in Terras's stopping
time distribution \cite{Terras1976} and appears in the survey literature
\cite{Lagarias1985,Korec1994,ApplegateLagarias1995}. What we claim as new is
the \emph{proof-complexity} statement: the transfer of the count from one
greedy tree to \emph{all} correct certificate systems (Theorem~A), the
two-sided non-asymptotic constants and the elementary proof of the exact
rate (Theorem~B), the Sturmian and automata-theoretic analysis of the
frontier (Theorems~C--E), the machine-verified core, and the verified
depth-$20$ artifact. In the language of proof complexity, Theorems~A--B are
a lower bound for a natural, restricted proof system for the $3x+1$
problem, in the spirit of lower bounds for resolution or bounded-depth
systems, and they rigorously prune a region of the strategy space in which
a large fraction of amateur and computer-assisted attempts operate.

\emph{Related recent work.} The parity-vector viewpoint has seen recent
activity close to our tools: Rozier and Terracol \cite{RozierTerracol2026}
study the \emph{paradoxical} orbit segments (affine coefficient below one,
yet no descent) which are exactly the exceptional integers of our collapse
lemma; Niu \cite{Niu2026} enumerates paradoxical windows in closed form and
observes their continued-fraction structure; and Fern\'andez and
Ib\'a\~nez \cite{FernandezIbanez2026} prove that Christoffel words are
extremal for a cycle functional. We indicate the points of contact as they
arise (Remarks~\ref{rem:countnum} and \ref{rem:sturmrelated}); the
proof-system lower bounds and the counting bound proved here do not appear
in these works.

The paper is organized as follows. Section~\ref{sec:prelim} fixes notation
and proves the Terras bijection in the form we need.
Section~\ref{sec:cert} defines certificates and systems.
Section~\ref{sec:climb} proves the core lemmas, Section~\ref{sec:counting}
the counting theorem, Section~\ref{sec:adaptive} its extension to
disjunctive and adaptive systems, and Section~\ref{sec:rate} the exact
growth rate.
Section~\ref{sec:sturmian} develops the Sturmian structure,
Section~\ref{sec:local} the local obstructions,
Section~\ref{sec:artifacts} describes the verified artifacts, and
Section~\ref{sec:discussion} delimits what is \emph{not} blocked by these
results. An appendix reproduces, for completeness, the standard cycle
bound in the form used here.

\section{Preliminaries: the compressed map and parity vectors}\label{sec:prelim}

Throughout, $T$ denotes the \emph{compressed} Collatz map on the positive
integers,
\[
T(n)=\begin{cases} n/2, & n \text{ even},\\[2pt] (3n+1)/2, & n \text{ odd},\end{cases}
\]
introduced by Terras \cite{Terras1976}. For odd $n$ one has $C^2(n)=T(n)$,
and since $3n+1>n$, the orbit of $n$ under $C$ descends below $n$ if and
only if the orbit under $T$ does; all our statements therefore transfer to
the standard map by this time change. We write
$\theta=\log_3 2 = 0.6309297536\ldots$ and
$H(x)=-x\log_2 x-(1-x)\log_2(1-x)$ for the binary entropy.

For $n\ge 1$ and $j\ge 0$ let $b_j(n)=T^j(n)\bmod 2$ be the \emph{parity
vector} of the orbit, and $e_j(n)=b_0(n)+\dots+b_{j-1}(n)$ the number of odd
steps among the first $j$.

\begin{lemma}[affine form and Terras bijection]\label{lem:terras}
For every $k\ge 0$ and every residue $r$ mod $2^k$:
\begin{enumerate}
\item[(a)] the parities $b_0,\dots,b_{k-1}$ are constant on the class
  $\{n\equiv r \pmod{2^k}\}$; write $e_j=e_j(r)$ for the resulting prefix
  counts;
\item[(b)] there are integers $\beta_j\ge 0$ ($0\le j\le k$), depending only
  on $r$, with
  \[
  T^j(n)=\frac{3^{e_j} n+\beta_j}{2^j} \qquad (n\equiv r \bmod 2^k,\ 0\le j\le k),
  \]
  where $\beta_0=0$ and $\beta_{j+1}=3\beta_j+2^j$ if $b_j=1$,
  $\beta_{j+1}=\beta_j$ if $b_j=0$; moreover $\beta_j\le 3^j$;
\item[(c)] the map $r\mapsto (b_0,\dots,b_{k-1})$ is a bijection from
  $\Z/2^k\Z$ onto $\{0,1\}^k$.
\end{enumerate}
\end{lemma}

\begin{proof}
Induction on $k$, all three claims jointly. For $k=0$ there is nothing to
prove. Assume the claims for $k$ and fix a class $r$ mod $2^k$, with affine
form $T^k(n)=(3^{e_k}n+\beta_k)/2^k$. If $n$ and $n'=n+2^k$ are the two
lifts of $r$ mod $2^{k+1}$, then
$T^k(n')-T^k(n)=3^{e_k}\cdot 2^k/2^k=3^{e_k}$, which is odd. Hence
$b_k$ is constant on each class mod $2^{k+1}$ and takes different values on
the two lifts of each class mod $2^k$; this proves (a) and (c) at level
$k+1$. For (b), if $b_k=0$ then $T^{k+1}(n)=T^k(n)/2=(3^{e_k}n+\beta_k)/2^{k+1}$;
if $b_k=1$ then
$T^{k+1}(n)=(3\,T^k(n)+1)/2=(3^{e_k+1}n+3\beta_k+2^k)/2^{k+1}$.
The update rules and $\beta_{j+1}\le 3\beta_j+2^j\le 3\cdot 3^j$ give
$\beta_j\le 3^j$ by induction.
\end{proof}

\section{Certificates and systems}\label{sec:cert}

\begin{definition}[dyadic affine descent certificate]\label{def:cert}
A \emph{certificate} is a tuple $(r,k,e,\beta,E)$ where $k\ge 1$, $r$ is a
residue mod $2^k$ whose parity vector has $e=e_k(r)$ odd steps and affine
constant $\beta=\beta_k(r)$ as in Lemma~\ref{lem:terras}, $E$ is a finite
set, and the \emph{descent condition}
\[
T^k(n)<n \qquad\text{for all } n\in \mathcal{C},\qquad
\mathcal{C}=\{n\ge 1: n\equiv r \bmod 2^k\}\setminus E ,
\]
holds. More generally we allow the certificate to be restricted to a class
$\mathcal{C}'\subseteq\mathcal{C}$ defined by an additional congruence to an
odd modulus. In all cases the \emph{dyadic atom} of the certificate is the
class $r$ mod $2^k$ in which it lives, and $k$ is its \emph{dyadic depth}.

A \emph{certificate system} $S$ is a finite or infinite family of
certificates, together with a finite set $D$ of directly verified integers,
such that the union of the certificate classes and $D$ covers
$\N\setminus\{1,2\}$.
\end{definition}

A certificate system is exactly what is needed to run the folk proof: if
every $n\ge 3$ lies in some certified class (or in $D$, whose orbits are
checked to reach $1$ directly), strong induction on $n$ proves the
conjecture, provided the exceptional sets are also checked. Whether a
correct system exists at all is therefore equivalent to the conjecture
itself (if the conjecture holds, the canonical first-descent tree is one);
the point of this paper is that \emph{any} correct system must be infinite,
with a quantified, exponentially growing frontier at every depth. The
verified artifact of Section~\ref{sec:artifacts} is a finite \emph{partial}
system: it certifies all classes mod $2^{20}$ except the $c(20)$ climbing
ones.

\section{Climbing classes and the core lemmas}\label{sec:climb}

\begin{definition}\label{def:climbing}
A residue $r$ mod $2^k$ (equivalently, by Lemma~\ref{lem:terras}(c), its
parity vector) is \emph{climbing} if $3^{e_j}>2^j$ for every
$1\le j\le k$; equivalently $e_j>\theta j$ for all $j\le k$. Let $c(k)$
denote the number of climbing residues mod $2^k$, and let
$L\subseteq\{0,1\}^*$ be the language of climbing words.
\end{definition}

The first lemma shows that climbing classes exist in abundance --- indeed
the class $-1$ mod $2^m$ climbs for $m-1$ full steps, realizing the
extremal growth rate $(3/2)^{m-1}$.

\begin{lemma}[climb lemma; machine-verified]\label{lem:climb}
Let $m\ge 1$, $t\ge 1$ and $n=2^m t-1$. Define $n_j=3^j 2^{m-j}t-1$ for
$0\le j\le m-1$, so $n_0=n$. Then for $0\le j\le m-2$: $n_j$ is odd,
$T(n_j)=n_{j+1}$, and $n_{j+1}>n_j$. In particular the first $m-1$ steps of
the orbit of $n$ are all odd steps, the orbit is strictly increasing along
them, and $n_j\equiv -1 \pmod{2^{m-j}}$ for all $j\le m-1$.
\end{lemma}

\begin{proof}
For $j\le m-1$ the number $3^j2^{m-j}t$ is even (as $m-j\ge 1$), so $n_j$
is odd. For $j\le m-2$,
\[
T(n_j)=\frac{3n_j+1}{2}=\frac{3^{j+1}2^{m-j}t-2}{2}=3^{j+1}2^{m-j-1}t-1=n_{j+1},
\]
and $T(n)=(3n+1)/2>n$ holds for every positive $n$. The congruence is read
off the formula. This lemma is machine-verified in Lean~4
(Section~\ref{sec:artifacts}).
\end{proof}

\begin{corollary}\label{cor:cknonzero}
For every $k\ge 1$ the class $-1$ mod $2^k$ is climbing; hence $c(k)\ge 1$.
\end{corollary}

\begin{proof}
By Lemma~\ref{lem:climb} (with $m=k$), the values $n_0,\dots,n_{k-1}$ are
all odd, so the parity vector of the class is all-ones: $e_j=j$ and
$3^j>2^j$ for every $j\le k$.
\end{proof}

\begin{lemma}[climbing classes admit no certificate]\label{lem:nocert}
If $3^e\ge 2^k$ and $\beta\ge 0$, then $(3^e n+\beta)/2^k<n$ holds for no
$n\ge 1$. Consequently a certificate of depth $k$ (which certifies descent
for all but finitely many $n$ in an infinite class) must have $3^{e_k}<2^k$,
i.e.\ its dyadic atom is non-climbing at its own depth.
\end{lemma}

\begin{proof}
$(3^e n+\beta)/2^k \ge n\cdot 3^e/2^k \ge n$.
\end{proof}

\begin{lemma}[prefix closedness and locality]\label{lem:locality}
\begin{enumerate}
\item[(a)] Climbing is prefix-closed: if $r$ mod $2^k$ is climbing and
  $k'\le k$, then the reduction $r$ mod $2^{k'}$ is climbing.
\item[(b)] A certificate of dyadic depth $k'\le k$ covers no element of any
  climbing class mod $2^k$.
\item[(c)] A certificate of dyadic depth $k''>k$ has its atom contained in
  exactly one class mod $2^k$.
\end{enumerate}
\end{lemma}

\begin{proof}
(a) is immediate from Definition~\ref{def:climbing}, since the conditions
for $j\le k'$ are a subset of those for $j\le k$ and, by
Lemma~\ref{lem:terras}, the parity prefix of the reduced class is the
prefix of the parity vector. For (b): by Lemma~\ref{lem:nocert} the
certificate's atom $r'$ mod $2^{k'}$ satisfies $3^{e_{k'}}<2^{k'}$, so it is
not climbing at depth $k'$; but by (a) the reduction mod $2^{k'}$ of any
climbing class mod $2^k$ \emph{is} climbing at depth $k'$; hence the two
classes mod $2^{k'}$ are distinct and therefore disjoint. (c) is trivial.
\end{proof}

\section{The counting theorem}\label{sec:counting}

\begin{thmA}
Let $S$ be any correct certificate system. Then for every $k\ge 1$ the
certificates of $S$ of dyadic depth exceeding $k$ have atoms lying in
every climbing class mod $2^k$ --- hence in at least $c(k)$ distinct
classes mod $2^k$. Consequently:
\begin{enumerate}
\item[(a)] any family of certificates all of dyadic depth $\le K$ leaves
  every climbing class mod $2^K$ uncovered, a set of natural density
  $c(K)/2^K$; in particular no finite certificate system exists;
\item[(b)] with Theorem~B below, the frontier of every correct system at
  depth $k$ has width $c(k)=2^{H(\theta)k+o(k)}$, and covering all but
  density $\varepsilon$ requires depth
  $K\ge \log_2(1/\varepsilon)/(1-H(\theta)+o(1))$.
\end{enumerate}
We stress that the count is a statement about \emph{distinct classes}:
in a correct system the number of certificates deeper than $k$ is
infinite for every $k$ (a finite number would bound all depths, which
(a) forbids), so the content of the theorem is the width of the frontier,
not a cardinality.
\end{thmA}

\begin{proof}
Fix $k$ and a climbing class $\mathcal{K}$ mod $2^k$. The class contains
infinitely many integers and only finitely many lie in the directly
verified set $D$, so some $n\in\mathcal{K}$ is covered by a certificate of
$S$. By Lemma~\ref{lem:locality}(b), no certificate of depth $\le k$ covers
any element of $\mathcal{K}$; hence that certificate has depth $>k$. By
Lemma~\ref{lem:locality}(c), a certificate of depth $>k$ meets exactly one
class mod $2^k$, so certificates assigned in this way to distinct climbing
classes are distinct. There are $c(k)$ climbing classes, each met by its own certificate of
depth $>k$. Claim (a): by Lemma~\ref{lem:locality}(b) no certificate of
depth $\le K$ covers any element of a climbing class mod $2^K$, and the
climbing classes have density $c(K)/2^K\ge 2^{-K}$
(Corollary~\ref{cor:cknonzero}); a finite system has a maximal depth.
Claim (b) follows from Theorem~B and Corollary~\ref{cor:density}.
\end{proof}

\begin{remark}\label{rem:thmA}
Theorem~A is machine-verified in an abstract certificate-system
formulation, with the count rendered as an injection from climbing
classes into deep certificates (Section~\ref{sec:artifacts}).
Sharpe's no-go theorem \cite{Sharpe2026} excludes, by the same all-ones
mechanism (the integers $2^L-1$ shadow the $2$-adic fixed point $-1$),
every scheme whose potential is $n\,W(n)$ with $W$ bounded above and
below and whose descent window is uniformly bounded --- the regime of a
\emph{finite} computation. Theorem~A addresses the complementary regime
of unbounded depth and quantifies it: the frontier width
$c(k)=2^{H(\theta)k+o(k)}$, its robustness under disjunction and adaptive
stopping (Theorem~A$'$), and its Sturmian and non-regular structure
(Theorems~C--E) have no counterpart in the bounded-window statement.
The proof uses nothing about the internal structure of the certificates
beyond Lemmas~\ref{lem:nocert} and \ref{lem:locality}; in particular it
applies verbatim to systems mixing arbitrary depths, refining by odd
moduli, and discarding arbitrary finite exceptional sets. Systems whose
members certify descent for a \emph{disjunction} of congruence classes
under an \emph{adaptive} stopping rule are not covered by the statement as
given; Section~\ref{sec:adaptive} extends the bound to exactly that
setting.
\end{remark}

\section{Robustness: disjunctions and adaptive stopping}\label{sec:adaptive}

The folk strategy admits two natural liberalizations that
Definition~\ref{def:cert} does not cover: a single certificate might claim
descent for a \emph{union} of congruence classes, and it might allow the
descent time to \emph{vary with $n$} instead of being the fixed depth of an
affine identity. This section shows that neither liberalization, nor both
together, weakens the counting bound: adaptive stopping collapses to a
canonical stopping time on every congruence class, and each constituent
class of a valid disjunction can serve at most one climbing class.

\begin{definition}[adaptive disjunctive certificate]\label{def:adaptive}
An \emph{adaptive disjunctive certificate} is a tuple $(U,K,E)$ where
$U=C_1\cup\dots\cup C_w$ is a finite union of congruence classes
$C_t=\{n\ge 1: n\equiv r_t \bmod 2^{d_t}q_t\}$ with $q_t$ odd
(we call $d_t$ the \emph{dyadic depth} of the constituent $C_t$, and $w$
the \emph{width}), $K\ge 1$ is the \emph{horizon}, $E$ is finite, and the
claim
\[
\forall\, n\in U\setminus E\ \ \exists\, j\le K:\ T^j(n)<n
\]
is true. An \emph{adaptive system} is a countable family of such
certificates, plus a finite directly verified set $D$, with
$\bigcup_i (U_i\setminus E_i)\cup D\supseteq\N\setminus\{1,2\}$.
Since every eventually periodic subset of $\N$ is, up to a finite set, a
finite union of congruence classes, this notion covers every certificate
whose scope is defined by congruence information of any modulus ---
including odd moduli such as powers of~$3$.
\end{definition}

\begin{lemma}[adaptive collapse]\label{lem:collapse}
Fix a class $r$ mod $2^K$ with prefix counts $e_j$
(Lemma~\ref{lem:terras}) and let
$j_1=\min\{j\le K: 3^{e_j}<2^j\}$ (with $j_1=\infty$ if the class is
climbing at depth $K$). Then for every $n\equiv r \pmod{2^K}$:
\begin{enumerate}
\item[(a)] $T^j(n)>n$ for every $j\le K$ with $3^{e_j}>2^j$ --- with no
  condition on $n$; this holds in particular for every $j<j_1$, and if
  $j_1=\infty$ then no $n\ge 1$ in the class descends below itself within
  $K$ steps;
\item[(b)] if $j_1\le K$ and $n>3^{j_1}$ then $T^{j_1}(n)<n$.
\end{enumerate}
Consequently, for $n>3^K$ the events ``$T^j(n)<n$ for some $j\le K$'' and
``$j_1\le K$'' coincide, and when $j_1\le K$ the first descent time equals
$j_1$:
on every congruence class, an adaptive stopping time collapses to the
canonical first-descent time. (This is the elementary, horizon-bounded
direction of Terras's coefficient stopping time analysis
\cite{Terras1976}.)
\end{lemma}

\begin{proof}
(a) $T^j(n)=(3^{e_j}n+\beta_j)/2^j\ge 3^{e_j}n/2^j>n$ whenever
$3^{e_j}>2^j$, with no condition on $n$; for $j<j_1$ this inequality holds
by minimality of $j_1$. (b) Descent at $j_1$ is
$n(2^{j_1}-3^{e_{j_1}})>\beta_{j_1}$; the left side is $\ge n$ (the bracket
is a positive integer) and $\beta_{j_1}\le 3^{j_1}$ by
Lemma~\ref{lem:terras}(b). The consequence follows since steps $j<j_1$
never descend by (a) and step $j_1$ descends by (b). Both halves, together
with the affine identity and the remainder bound $\beta_j+2^j\le 3^j$, are
machine-verified in Lean~4 (Section~\ref{sec:artifacts}).
\end{proof}

\begin{lemma}[persistence of climbing]\label{lem:persist}
If $r$ mod $2^k$ is climbing, its extension appending the parity $1$ is a
climbing class mod $2^{k+1}$. Hence $c(k+1)\ge c(k)$; every climbing class
mod $2^k$ contains a climbing class mod $2^K$ for every $K\ge k$; and the
intersection of a climbing class with any congruence class to an odd
modulus is infinite.
\end{lemma}

\begin{proof}
Prefix conditions up to $k$ are unchanged, and at $k+1$:
$3^{e_k+1}=3\cdot 3^{e_k}>3\cdot 2^k>2^{k+1}$. Iterating gives the second
claim; the third is the Chinese remainder theorem. The arithmetic step is
machine-verified (Section~\ref{sec:artifacts}).
\end{proof}

\begin{thmAp}
Let $S$ be any correct adaptive system. Then:
\begin{enumerate}
\item[(a)] every constituent class of every certificate of $S$ is
  non-climbing at its own dyadic depth;
\item[(b)] for every $k\ge 1$, the constituent classes of dyadic depth
  exceeding $k$ occupy every climbing class mod $2^k$ (hence at least
  $c(k)$ distinct classes), and every certificate covering an element of a
  climbing class mod $2^k$ has horizon exceeding $k$;
\item[(c)] (packaging invariance) for every valid $(U,K,E)$,
  \[
  U\setminus\bigl(E\cup[1,3^K]\bigr)\ \subseteq\ D(K):=
  \bigcup\{\,\text{classes mod } 2^K \text{ with } j_1\le K\,\},
  \]
  and conversely $D(K)$ is covered by the canonical uniform tree of
  first-descent certificates of depth $\le K$ with exceptions in
  $[1,3^K]$. Every adaptive disjunctive system is therefore, up to finite
  sets, a repackaging of the canonical first-descent tree.
\end{enumerate}
\end{thmAp}

\begin{proof}
(a) Suppose the constituent $C_t$ (dyadic depth $d$, odd modulus $q_t$,
inside a certificate of horizon $K$) is climbing at depth $d$. By
Lemma~\ref{lem:persist} there is a climbing class mod $2^{\max(d,K)}$
inside its dyadic class, and its intersection with the congruence
condition mod $q_t$ is infinite; every $n$ in that intersection satisfies
$T^j(n)>n$ for all $j\le K$ by Lemma~\ref{lem:collapse}(a). All but
finitely many of these $n$ avoid $E$, so the certificate's claim is false.

(b) Fix a climbing class $\mathcal{K}$ mod $2^k$ and pick
$n\in\mathcal{K}\setminus D$ covered by some certificate, say
$n\in C_t\setminus E$. By (a) the class $C_t$ has a first non-climbing
prefix $j_1\le d_t$. Since $n\in\mathcal{K}$, the parity prefix of $n$ is
climbing at every depth $\le k$; since $n\in C_t$, its prefix at depth
$j_1$ is that of $C_t$, which is non-climbing. Hence $j_1>k$, so
$d_t\ge j_1>k$, and the dyadic class of $C_t$ reduces mod $2^k$ to the
single class $\mathcal{K}$. Distinct climbing classes therefore meet
distinct constituent classes, so the constituents of dyadic depth $>k$
occupy every climbing class mod $2^k$. For the horizon: the prefixes of $n$ at depths
$\le k$ are climbing, so $T^j(n)>n$ for all $j\le k$
(Lemma~\ref{lem:collapse}(a)), and the certified descent time of $n$
exceeds $k$; hence $K>k$.

(c) For $n\in U\setminus E$ with $n>3^K$, the claim provides a descent
within $K$ steps, and Lemma~\ref{lem:collapse} shows this forces
$j_1(n\bmod 2^K)\le K$, i.e.\ $n\in D(K)$. Conversely each class mod
$2^K$ with $j_1\le K$ is certified by the affine certificate of depth
$j_1$ at its first-descent prefix, whose exceptional set is contained in
$[1,3^{j_1}]$ by Lemma~\ref{lem:collapse}(b).
\end{proof}

\begin{remark}[numerical verification]\label{rem:adaptivenum}
The collapse of Lemma~\ref{lem:collapse} was checked exhaustively for all
$2\le n\le 2^{20}$ at horizon $K=12$ (no violations; in this range no
exceptional stragglers below the $3^{j_1}$ bound occurred at all), the
monotonicity $c(k+1)\ge c(k)$ was confirmed to $k=20$ with the anchor
$c(20)=27328$ of Section~\ref{sec:artifacts}, and invalidity witnesses for
climbing constituents with odd moduli $q\in\{5,7,11\}$ were confirmed at
horizons up to $80$.
\end{remark}

The collapse lemma has a further unconditional consequence: a power-saving
upper bound on the number of potential counterexamples to the conjecture
itself.

\begin{corollary}[counting potential counterexamples]\label{cor:count}
Let $\mathcal{E}=\{n>1:\ T^j(n)\ge n \text{ for all } j\ge 1\}$, the set of
integers whose orbit never descends below its starting point. (The least
element of any divergent orbit and of any nontrivial cycle lies in
$\mathcal{E}$, and the conjecture is equivalent to
$\mathcal{E}=\emptyset$.) Then, unconditionally and with an absolute constant $C$,
\[
|\mathcal{E}\cap[1,N]|\ \le\ C\,N^{\,H(\theta)},\qquad
H(\theta)=0.9499555\ldots :
\]
the counting exponent equals the certificate growth rate $\gamma$.
\end{corollary}

\begin{proof}
Let $n\in\mathcal{E}\cap(M,2M]$ with $M=2^m$, and let $e$ be the number
of odd steps among the first $m$ steps of the orbit of $n$. Every iterate
$n'$ is $\ge n>M$, so an odd step multiplies the current value by
$(3n'+1)/(2n')\le (3+1/M)/2$ and an even step by $1/2$; hence
\[
1\ \le\ \frac{T^m(n)}{n}\ \le\ 2^{-m}\Bigl(3+\frac1M\Bigr)^{e},
\qquad\text{i.e.}\qquad e\ \ge\ m\,\theta_M,\quad
\theta_M:=\frac{\log 2}{\log(3+1/M)}=\theta-O(1/M).
\]
By Lemma~\ref{lem:terras}(c), each parity vector of length $m$ occurs
for exactly one integer of $(M,2M]$, so
$|\mathcal{E}\cap(M,2M]|\le\sum_{e\ge m\theta_M}\binom{m}{e}
\le 2^{mH(\theta_M)}$ (the entropy bound, valid since $\theta_M>1/2$).
As $m\,(H(\theta_M)-H(\theta))=O(m/M)\to 0$, this is at most
$C_0\,2^{mH(\theta)}=C_0M^{H(\theta)}$ with an absolute $C_0$, and summing
over the dyadic blocks $M=2^m\le N$ gives the claim.
\end{proof}

\begin{remark}\label{rem:countnum}
(i) The congruence-shadow route gives the same exponent with more effort:
Lemma~\ref{lem:collapse} forces every $n\in\mathcal{E}$ to be climbing at
all depths $j\le\log_3 n$, which yields the weaker exponent
$1-\theta(1-H(\theta))=0.9684254\ldots$; forcing climbing up to depth
$\log_2 n$ (and hence $H(\theta)$ again) requires bounding the
\emph{margin} $B(j):=\max\lfloor\beta_j/(2^j-3^{e_j})\rfloor$ over classes
whose first non-climbing depth is $j$, e.g.\ through the effective
irrationality measure $|j\log 2-e\log 3|>c\,j^{-13.3}$ of Rhin
\cite{Rhin1987}. The integers that dodge descent at their first
non-climbing depth are initial terms of \emph{paradoxical sequences} in
the sense of Rozier--Terracol \cite{RozierTerracol2026} --- a proper
subclass, since a paradoxical window may also end after a descent (for
$n=7$ the coefficient first drops below one at $j=7$, where $T^7(7)=5<7$,
yet the length-$8$ window with $T^8(7)=8$ is paradoxical). Their
remainder bound \cite[Thm.~2.4]{RozierTerracol2026},
$\beta_j/2^j\le(3^{q}-2^{q})/2^{q}$ with $q=e_j$, bounds the margin by
$(3/2)^{\theta j}j^{13.3}$ up to constants, and their Theorem~1.3 shows
no paradoxical sequence starts in $(4614,\,2.8\times 10^{19})$.
Numerically $B(j)\le 24$ for all $j\le 24$ (the real-valued maximum of
$\beta_j/(2^j-3^{e_j})$ is $24.54$ at $j=8$), with spikes at the depths
where $2^j$ is unusually close to a power of $3$ from above (e.g.\
$j=8$: $2^8/3^5=256/243$, the Pythagorean limma) --- the same
continued-fraction structure observed for paradoxical ratios in
\cite{Niu2026}. (ii) The argument of the corollary is Terras's \cite{Terras1976}: the
parity vector of length $m$ is a bijection on $(M,2M]$, and non-descent
forces a supercritical density of odd steps. Terras's theorem in this
finite form was recently machine-verified by Sharpe \cite{Sharpe2026},
with the explicit bound
$\#\{n\in[2^k,2^{k+1}):\text{no descent within } k \text{ steps}\}
\le 3^k/2^{\lfloor 17(k-1)/27\rfloor}\approx 2^{0.955k}$, where the
rational approximation $3^{17}\le 2^{27}$ replaces the entropy function;
the corollary sharpens the exponent to the exact entropy value
$H(\theta)=0.94996\ldots$ with an absolute constant and identifies it
with the certificate growth rate $\gamma$. The recent parity-vector
literature enumerates paradoxical \emph{windows} \cite{Niu2026} and shows
that finiteness of the set of paradoxical sequences implies the
conjecture \cite[Cor.~3.3]{RozierTerracol2026}. We are not aware of a
stated count of $\mathcal{E}$ with the exponent $H(\theta)$, but the
bound is a routine consequence of Terras's method, and we state it with
that caveat; its combinatorial core is machine-verified
(Section~\ref{sec:artifacts}). (iii) The shadow
method is exhausted at the exponent $H(\theta)$:
the set of integers in $(M,2M]$ satisfying all depth-$m$ congruence
necessary conditions has exactly $c(m)$ elements, so no argument using
congruence shadows alone can lower the exponent further. Pushing the
forced-climbing depth to its true limit $\log_{3/2}n$ would require
counting climbing residues that are \emph{small as integers} ---
information not contained in shadows (heuristically that limit exponent is
$1-(1-H(\theta))/\log_2(3/2)=0.91445\ldots$). (iv) Thematically, the corollary ties the two halves of the problem
together: the count of potential counterexamples is governed by the same
entropy that prices the certificate frontier, and the finer margin
analysis of (i) runs through the same linear-forms data that governs the
cycle bound of Appendix~\ref{app:cycle}.
\end{remark}

\begin{remark}[what robustness does and does not say]\label{rem:adaptivescope}
Theorem~A$'$ closes every liberalization of the folk strategy whose
certificate scopes are eventually periodic sets --- in particular,
partitioning by residues modulo powers of $3$, or any mixed modulus, buys
nothing, and neither does any orbit-adaptive stopping rule. What remains
outside its quantification are certificates whose scopes are \emph{not}
unions of congruence classes: sets defined by inequalities on orbit
values, by unbounded-window parity information, or by genuinely analytic
(ensemble) data. Section~\ref{sec:discussion} discusses these exits.
\end{remark}

\section{The exact growth rate}\label{sec:rate}

\begin{lemma}[rotation lemma]\label{lem:rotation}
Let $b=(b_1,\dots,b_k)\in\{0,1\}^k$ with $e=\sum_i b_i$ and $e-\theta k>0$.
Then some cyclic rotation of $b$ is climbing, i.e.\ has all prefix sums
$e_j-\theta j$ positive.
\end{lemma}

\begin{proof}
Set $w_i=b_i-\theta$ and $S_j=w_1+\dots+w_j$, $S_0=0$, so
$S_k=e-\theta k>0$. Since $\theta$ is irrational, the values
$S_0,\dots,S_{k-1}$ are pairwise distinct (an equality $S_a=S_b$ with
$a<b$ would make $\theta=(e_b-e_a)/(b-a)$ rational). Let $j^*$ be the
unique index in $\{0,\dots,k-1\}$ at which $\min_{0\le j\le k-1}S_j$ is
attained, and consider the rotation
$b'=(b_{j^*+1},\dots,b_k,b_1,\dots,b_{j^*})$. Its prefix sums are
\[
S'_i=S_{j^*+i}-S_{j^*}\quad (1\le i\le k-j^*),\qquad
S'_{k-j^*+i}=S_k-S_{j^*}+S_i \quad (1\le i\le j^*).
\]
The first group is positive because $S_{j^*}$ is the strict minimum among
$S_0,\dots,S_{k-1}$ and $S_k>S_0\ge S_{j^*}$; the second group is positive
because $S_i\ge S_{j^*}$ and $S_k>0$. Hence all prefix sums of $b'$ are
positive, which is the climbing condition.
\end{proof}

\begin{remark}[machine verification; irrationality is removable]
\label{rem:rotlean}
The lemma is machine-verified (Section~\ref{sec:artifacts}) in a purely
integer form: comparisons $S_a\le S_b$ are encoded as the cross-multiplied
inequalities $3^{e_a}2^{b}\le 3^{e_b}2^{a}$, and the rotated word's prefix
conditions become the two families of integer inequalities produced by the
proof. The formalization selects the \emph{last} index attaining the
minimum, which yields the required strictness on the first group directly
--- so the appeal to the irrationality of $\theta$ in the proof above is a
convenience of the real-variable formulation, not a genuine hypothesis.
\end{remark}

\begin{thmB}
For every $k\ge 1$, with $\theta=\log_3 2$,
\[
\frac1k\binom{k}{\ceil{\theta k}} \;\le\; c(k)\;\le\;(k+1)\binom{k}{\ceil{\theta k}} .
\]
Consequently $\gamma:=\lim_{k\to\infty}\frac1k\log_2 c(k)$ exists and
\[
\gamma=H(\theta)=H(\log_3 2)=0.9499555272\ldots
\]
Moreover $c$ is supermultiplicative, $c(k_1+k_2)\ge c(k_1)c(k_2)$, so the
limit also equals $\sup_k \frac1k \log_2 c(k)$.
\end{thmB}

\begin{proof}
\emph{Lower bound.} Let $e=\ceil{\theta k}$, so $e-\theta k>0$ ($\theta k$
is irrational). Partition the $\binom{k}{e}$ words with $e$ ones into
classes under cyclic rotation; each class has at most $k$ elements, and by
Lemma~\ref{lem:rotation} each class contains at least one climbing word.
By Lemma~\ref{lem:terras}(c) every word is the parity vector of exactly one
residue, so $c(k)\ge \binom{k}{e}/k$.

\emph{Upper bound.} A climbing word satisfies $e_k>\theta k$, so
$e_k\ge\ceil{\theta k}$. Since $\theta>1/2$, the binomial coefficients
$\binom{k}{e}$ are non-increasing for $e\ge\ceil{\theta k}\ge k/2$, whence
\[
c(k)\le \sum_{e\ge \ceil{\theta k}}\binom{k}{e}
\le (k+1-\ceil{\theta k})\binom{k}{\ceil{\theta k}}\le (k+1)\binom{k}{\ceil{\theta k}}.
\]

\emph{Supermultiplicativity.} If $u$ is climbing of length $k_1$ and $v$
climbing of length $k_2$, then the concatenation $uv$ is climbing: for
$j\le k_1$ the prefix sums are those of $u$; for $j=k_1+i$ they equal
$S^u_{k_1}+S^v_i>0$. Distinct pairs give distinct concatenations, so
$c(k_1+k_2)\ge c(k_1)c(k_2)$. Fekete's lemma applied to the superadditive
sequence $\log_2 c(k)$ (finite since $c(k)\le 2^k$, and $c(k)\ge 1$ by
Corollary~\ref{cor:cknonzero}) yields
$\lim_k \frac1k\log_2 c(k)=\sup_k\frac1k\log_2 c(k)$.

\emph{The value.} By Stirling,
$\frac1k\log_2\binom{k}{\ceil{\theta k}}\to H(\theta)$, and the polynomial
factors $1/k$, $(k+1)$ do not affect the limit.
\end{proof}

\begin{remark}[numerical verification]\label{rem:numB}
The sandwich was verified computationally for all $k\le 1600$ against exact
values of $c(k)$ computed by dynamic programming over the walk of
Theorem~C, and the climbing condition was checked against exhaustive
enumeration of all $2^k$ words for $k\le 16$. The fluctuation
$\tilde b(k)=\log_2 c(k)-\log_2\binom{k}{\ceil{\theta k}}$ stays within the
proven window $[-\log_2 k,\ \log_2(k+1)]$ throughout.
\end{remark}

\begin{corollary}[density; classical in content]\label{cor:density}
The natural density of the union of climbing classes at depth $k$ is
$c(k)/2^k=2^{-(1-H(\theta))k+o(k)}$, with $1-H(\theta)=0.0500444728\ldots$;
equivalently, the greedy certificate tree of depth $k$ certifies descent on
a set of density $1-c(k)/2^k\to 1$. This recovers, in certificate form, the
density theorems of Terras and Everett
\cite{Terras1976,Everett1977,Korec1994}.
\end{corollary}

\section{Sturmian structure of the frontier}\label{sec:sturmian}

For $j\ge 1$ let
\[
\sigma_j=\floor{(j+1)\theta}-\floor{j\theta}\in\{0,1\}
\]
be the $j$-th letter of the \emph{characteristic Sturmian word} of slope
$\theta$ \cite[Ch.~2]{Lothaire2002}.

\begin{thmC}
For a word $(b_1,\dots,b_k)$ define the \emph{excess walk}
$d_j=e_j-\floor{\theta j}-1$. Then the word is climbing if and only if
$d_j\ge 0$ for all $1\le j\le k$, and the walk updates by
\[
d_{j+1}-d_j=b_{j+1}-\sigma_j\in
\begin{cases}
\{0,+1\}, & \sigma_j=0,\\
\{-1,0\}, & \sigma_j=1 .
\end{cases}
\]
Hence $c(k)$ equals the number of non-negative paths of the
$\sigma$-driven walk started from the forced first step $b_1=1$
(i.e.\ $d_1=0$).
\end{thmC}

\begin{proof}
Since $\theta j$ is irrational for $j \geq 1$, $e_j>\theta j$ is equivalent to
$e_j\ge\floor{\theta j}+1$, i.e.\ $d_j\ge 0$. The update is immediate from
$e_{j+1}=e_j+b_{j+1}$ and the definition of $\sigma_j$. At $j=1$:
$\floor{\theta}=0$, so $d_1=b_1-1\ge 0$ forces $b_1=1$.
\end{proof}

For the block structure, write the continued fraction
$\theta=[0;a_1,a_2,a_3,\dots]=[0;1,1,1,2,2,3,1,5,2,\dots]$ and let
$s_0=0$, $s_1=1$, $s_{n+1}=s_n^{\,a_{n+1}}s_{n-1}$ be the associated
standard words; their lengths are the continued-fraction denominators
$q_n=1,1,2,3,8,19,65,84,485,1054,\dots$ of $\theta$, and the
characteristic word $\sigma$ is their limit \cite[Ch.~2]{Lothaire2002}.
Encode the two letter steps of Theorem~C as infinite $0$--$1$ transfer
matrices acting on the state $d\in\{0,1,2,\dots\}$, in the row-vector
convention:
\[
M(0)_{d,d'}=[d'=d]+[d'=d+1],\qquad
M(1)_{d,d'}=[d'=d-1]_{d\ge1}+[d'=d] .
\]
Products of $k$ such matrices are banded (bandwidth $\le k$), so all
entries below are finite sums.

\begin{thmD}
Let $B_n=M(\sigma\text{-letters of } s_n)$ be the ordered product of letter
matrices along the standard word $s_n$ (row-vector convention: products
compose in word order). Then
\[
B_{n+1}=B_n^{\,a_{n+1}}B_{n-1}\qquad (n\ge 1),
\]
and for every $k$,
\[
c(k)=\Bigl(e_0^{\top} \,M(\sigma_1)M(\sigma_2)\cdots M(\sigma_{k-1})\Bigr)\mathbf 1 ,
\]
where $e_0$ is the coordinate vector of the state $d=0$.
\end{thmD}

\begin{proof}
The map assigning to a finite word the ordered product of its letter
matrices is a monoid morphism (in the row-vector convention, concatenation
maps to product in the same order), so the matrix recursion is the image of
the standard-word recursion $s_{n+1}=s_n^{\,a_{n+1}}s_{n-1}$. The formula
for $c(k)$ is Theorem~C: the matrix entry
$(M(\sigma_1)\cdots M(\sigma_{k-1}))_{0,d}$ counts the walks from $d_1=0$
to $d_k=d$ that stay non-negative, and summing over $d$ counts all
climbing words.
\end{proof}

\begin{remark}[numerics]\label{rem:sturmnum}
The identification of $\sigma$ with the limit of the standard words was
checked letter-by-letter to length $1054$ $(=q_9)$; the matrix recursion
was verified exactly on truncation-exact corners and modulo the prime
$3000017$ for $n\le 6$; and the matrix-product formula for $c(k)$ was
checked against the direct dynamic program for $k\le 200$. A two-letter
statistic captures most of the fluctuation of $\log_2 c(k)$ around
$H(\theta)k$: empirically
$\mathbb{E}[\Delta\varphi\mid\sigma=0]=+0.05004\approx 1-H(\theta)$ and
$\mathbb{E}[\Delta\varphi\mid\sigma=1]=-0.03022$ explain $82\%$ of the
variance, rising to $97.6\%$ at the level of derived-word blocks --- the
fluctuation of the certificate count is governed by the first two levels of
the continued-fraction hierarchy of $\log_2 3$.
\end{remark}

\begin{thmE}
The language $L$ of climbing words is not regular. More precisely, for
$a,b\ge 1$,
\[
1^a0^b\in L \iff b\le \Bigl\lceil a\,\tfrac{1-\theta}{\theta}\Bigr\rceil-1 ,
\]
and since the right-hand side is unbounded in $a$, the Myhill--Nerode
equivalence has infinitely many classes.
\end{thmE}

\begin{proof}
For the word $1^a0^b$ the prefix sums are $e_j-\theta j=j(1-\theta)>0$ for
$j\le a$, and $e_{a+i}-\theta(a+i)=a-\theta(a+i)$ for $1\le i\le b$, which
is positive iff $i<a(1-\theta)/\theta$; as $\theta$ is irrational, equality
is impossible, so the condition on $b$ is
$b\le\ceil{a(1-\theta)/\theta}-1$. Write $\phi(a)$ for this bound; $\phi$
is unbounded. If $a\ne a'$ with $\phi(a)<\phi(a')$, the word $0^{\phi(a')}$
separates $1^{a}$ from $1^{a'}$; hence infinitely many Myhill--Nerode
classes, and $L$ is not regular.
\end{proof}

\begin{remark}[related work]\label{rem:sturmrelated}
Balanced-word structures have independently surfaced on the cycle side of
the problem: Fern\'andez--Ib\'a\~nez \cite{FernandezIbanez2026} prove that
Christoffel words --- the finite relatives of standard Sturmian words ---
are, up to rotation, the unique maximizers of a cycle functional on binary
words of fixed length and weight, yielding bounds on the minimum element
of a hypothetical cycle. Their extremal role there and the driving role of
the characteristic Sturmian word here (Theorems~C--D) are distinct
phenomena with a common source: the affine functional $3^{e}/2^{j}$ is
extremized under balance constraints governed by the continued fraction of
$\log_2 3$. In the machine-verified setting, Sharpe \cite{Sharpe2026}
excludes specific mechanical (Sturmian) parity itineraries of slopes above
the critical line --- slope $1/\sqrt2$ and an explicit metallic family ---
as itineraries of actual orbits; that is a statement about orbits, whereas
Theorems~C--E concern the frontier of certificate systems, whose slope is
the critical $\theta$ itself.
\end{remark}

\begin{conjecture}\label{conj:precursive}
The sequence $c(k)$ is not P-recursive (holonomic). (Supporting evidence: an
exhaustive search found no polynomial-coefficient recurrence of order
$\le 6$ and degree $\le 6$, with the kernel computations confirmed modulo
three independent primes; and the Sturmian driving of Theorem~C ties the
fine structure of $c(k)$ to the continued fraction of $\log_2 3$, which is
not eventually periodic as far as it has been computed.)
\end{conjecture}

\section{Local obstructions}\label{sec:local}

The counting theorem bounds one specific proof system. It is natural to ask
whether small modifications escape it: a bounded ``syndrome correction'' to
the potential $\log_2 n$, or a hierarchy of finite-window constraints meant
to exclude the climbing statistics. The following two results close the
bounded-memory variants of both ideas; their scope limitations are stated
explicitly and discussed in Section~\ref{sec:discussion}.

\begin{thmF}
Let $d\ge 1$, $\psi:\Z/2^d\Z\to\R$, $\varepsilon>0$, and set
$\varphi(n)=\log_2 n+\psi(n\bmod 2^d)$. Then the one-step descent property
\[
\varphi(T(n))\le\varphi(n)-\varepsilon\qquad\text{for all } n>1
\]
fails. In fact it fails on the single class $n\equiv -1 \pmod{2^{d+1}}$:
there the syndrome is a fixed point and the step is expanding.
\end{thmF}

\begin{proof}
Let $n=2^{d+1}t-1$, $t\ge 1$. Then $n$ is odd, $n\equiv -1\pmod{2^d}$, and
by Lemma~\ref{lem:climb} (with $m=d+1$) $T(n)=3\cdot 2^{d}t-1\equiv -1
\pmod{2^d}$. Hence $\psi(T(n)\bmod 2^d)=\psi(n\bmod 2^d)$ and
\[
\varphi(T(n))-\varphi(n)=\log_2\frac{3n+1}{2n}>\log_2\frac32>0 ,
\]
contradicting descent by any $\varepsilon>0$.
\end{proof}

\begin{remark}[robustness of the obstruction]\label{rem:simplex}
Theorem~F is stated for the all-ones syndrome because one class suffices,
but the obstruction is much fatter: in the natural linear-programming
relaxation over the syndrome transition graph, every shift-invariant
measure on parity sequences with odd-step frequency $\rho>\theta$ has
positive drift $\rho\log_2 3-1>0$, and these measures form a non-empty
compact convex set. Any per-step, bounded-memory dyadic correction thus
fails robustly, not by a degenerate coincidence. By contrast, deleting the
states with two trailing ones makes the maximum mean cycle negative: an
explicit finite Lyapunov function exists for orbits that never visit
$n\equiv 3\pmod 4$; compare Corollary~\ref{cor:freq}.
\end{remark}

\begin{thmG}
Fix a window length $w\ge 1$ and a syndrome depth $d\ge 1$. There exist
infinitely many $n$ (namely all $n=2^{w+d+1}t-1$, $t\ge 1$) whose first $w$
compressed steps form a strictly increasing orbit segment with all steps
odd, unit $2$-adic valuations ($v_2(3n_j+1)=1$), and constant all-ones
syndrome $n_j\equiv -1 \pmod{2^d}$ for $0\le j\le w$. Consequently, any
family of constraints on length-$w$ windows of (parities, valuations,
depth-$d$ syndromes) that is valid for all non-descending integer orbit
segments is satisfied by the all-ones statistics: finite-window constraint
hierarchies cannot exclude climbing statistics.
\end{thmG}

\begin{proof}
Apply Lemma~\ref{lem:climb} with $m=w+d+1$: for $j\le w$ the values
$n_j=3^j2^{m-j}t-1$ satisfy $n_j\equiv -1\pmod{2^{m-j}}$ and $m-j\ge d+1$,
so all listed statistics of the window coincide with those of the all-ones
cylinder, while the segment is a genuine orbit of a positive integer and is
strictly increasing (hence non-descending).
\end{proof}

\begin{remark}[scope]\label{rem:scopeG}
Theorem~G is a $\forall w\,\exists n$ statement: for each window length
some integer realizes the climbing statistics. It does not produce a single
$n$ realizing them for all $w$ (that would be a divergent orbit). Its force
is contrapositive: any proof method whose reach is exhausted by constraints
on bounded windows of bounded-depth dyadic data cannot prove the
conjecture, since such constraints hold equally in the $2$-adic completion,
where the all-ones orbit $(-1\mapsto -1)$ exists. Local statistics see only
$\Z_2$, not $\N$.
\end{remark}

\begin{corollary}[forced visits to $3 \bmod 4$]\label{cor:freq}
If $n>3^k$ and $T^j(n)\ge n$ for all $j\le k$, then the orbit segment
visits $\{m\equiv 3 \pmod 4\}$ at least $(2\theta-1)k-1$ times among
$j=0,\dots,k-1$; the frequency is at least
$2\theta-1-O(1/k)$, with $2\theta-1=0.2618595\ldots$
\end{corollary}

\begin{proof}
Non-descent through depth $k$ forces the climbing condition $e_j>\theta j$
for all $j\le k$: if some prefix had $3^{e_j}<2^j$, then by
Lemma~\ref{lem:terras}(b) descent $T^j(n)<n$ would hold whenever
$n(2^j-3^{e_j})>\beta_j$, which is guaranteed by $n>3^k\ge\beta_j$
since $2^j-3^{e_j}\ge 1$. An odd $m$
has $T(m)$ odd iff $3m+1\equiv 2\pmod 4$ iff $m\equiv 3\pmod 4$. Thus
visits to $3\bmod 4$ within the window are exactly the adjacent pairs
$b_j=b_{j+1}=1$. A binary word of length $k$ with $e$ ones has at least
$e-(k-e+1)=2e-k-1$ such pairs (the ones form at most $k-e+1$ blocks). With
$e\ge\theta k$ this gives $(2\theta-1)k-1$.
\end{proof}

\section{Verified computational artifacts}\label{sec:artifacts}

\subsection{The depth-20 certificate system}
A greedy certificate tree of dyadic depth $20$ was generated and
independently verified. It consists of $4404$ certificates whose atoms,
together with the $c(20)=27328$ climbing classes mod $2^{20}$, form an
exact partition of $\Z/2^{20}\Z$; the certified classes have natural
density $1-27328/2^{20}=97.39\%$, and the single exceptional integer across
the entire system is $n=1$ (the endpoint of the trivial cycle, for which
descent is not required). Each certificate is stored as the full tuple
$(k,\ \mathrm{mod},\ r,\ e,\ \beta,\ \text{exceptions})$ in a JSON artifact.

The verifier is deliberately independent of the generator: it re-derives
the parity vectors and affine data symbolically from scratch
(Lemma~\ref{lem:terras}), checks exact cover and disjointness of the atoms,
re-checks the descent inequality $n(2^k-3^e)>\beta$ at the minimal element
of each class (which suffices for the whole class by monotonicity of the
affine form), and confirms the climbing condition for each unresolved
class. All arithmetic is exact (arbitrary-precision integers).

\subsection{Machine verification of the core lemmas}
The arithmetic cores of Theorems~A and A$'$ are verified in Lean~4,
compiled with core Lean only (no external libraries), toolchain pinned to
\texttt{leanprover/lean4:v4.32.2}, with no \texttt{sorry}:
\begin{itemize}
\item \texttt{Barrier.lean}: Lemma~\ref{lem:climb} as four theorems
  (\texttt{climb}, \texttt{climb\_mod}, \texttt{climb\_ge},
  \texttt{barrier}); axioms exactly \texttt{[propext, Quot.sound]}.
\item \texttt{Collapse.lean}: the affine identity of
  Lemma~\ref{lem:terras}(b) for the full compressed map
  (\texttt{affine}), the remainder bound $\beta_j+2^j\le 3^j$
  (\texttt{beta\_bound}), both halves of the collapse lemma
  (Lemma~\ref{lem:collapse}) as \texttt{climbing\_no\_descent} ($2^j<
  3^{e_j}\Rightarrow T^j(n)>n$ for every $n\ge 1$) and
  \texttt{descent\_forced} ($3^{e_j}<2^j$ and $n>3^j\Rightarrow
  T^j(n)<n$), and the arithmetic core of the persistence lemma
  (Lemma~\ref{lem:persist}) as \texttt{climb\_extend}; axioms within
  \texttt{[propext, Classical.choice, Quot.sound]}, the standard Lean
  axioms.
\item \texttt{Rotation.lean}: the rotation lemma
  (Lemma~\ref{lem:rotation}) in the integer cross-multiplied form of
  Remark~\ref{rem:rotlean} (\texttt{rotation}, built on a verified
  last-minimum scan \texttt{am} with its minimality and strictness
  lemmas); axioms within \texttt{[propext, Classical.choice,
  Quot.sound]}. The correspondence between the cross-multiplied
  conclusion and the literal climbing of the rotated word was
  additionally cross-checked exhaustively for all words of length
  $\le 14$.
\item \texttt{Counting.lean}: \emph{Theorem A itself}, in an abstract
  certificate-system formulation. A system is abstracted as an index
  type $I$ with depth $\mathrm{dep}\,i$, atom $\mathrm{atom}\,i$, and a
  coverage predicate, subject to exactly the three hypotheses used in
  the proof: containment in the dyadic atom, descent at the
  certificate's own depth on every covered integer, and coverage of all
  $n\ge N_0$ (which absorbs exceptional sets, odd-modulus refinements,
  and the finite directly verified set). The conclusion
  (\texttt{counting\_lower\_bound}) is constructive: an
  \emph{injection} from the climbing residues mod $2^K$ into
  $\{i: \mathrm{dep}\,i>K\}$ --- the statement ``at least $c(K)$
  certificates of depth $>K$'' rendered without cardinality machinery
  --- together with the corollary \texttt{no\_finite\_system} ($\forall
  K\,\exists i,\ \mathrm{dep}\,i>K$, witnessed by the all-ones class).
  The key lemma is \texttt{odds\_mod}: the first $j$ parities depend
  only on $n \bmod 2^j$ (the part of Lemma~\ref{lem:terras} the proof
  uses). Axioms within \texttt{[propext, Classical.choice,
  Quot.sound]}.
\item \texttt{TheoremB.lean} (separate Lake project, Lean
  \texttt{v4.33.1} with Mathlib pinned at the matching release):
  \emph{the two-sided sandwich of Theorem~B}. Words are
  \texttt{List.Vector Bool k}, the
  climbing count $c(k)$ is the cardinality of the climbing words, and
  $e_0=\min\{e: 2^k<3^e\}$; the theorem \texttt{theoremB} states
  $\binom{k}{e_0}\le k\cdot c(k)$ and $c(k)\le (k+1)\binom{k}{e_0}$ for
  $k\ge 1$. The lower bound composes the integer rotation theorem with
  the list identities for prefix counts of \texttt{List.rotate} and the
  fiber bound of the map ``word $\mapsto$ its climbing rotation''
  (\texttt{Finset.card\_le\_mul\_card\_image}); the upper bound uses
  $2e_0>k$ and the monotone decrease of binomial coefficients above the
  middle; the identification of $\#\{\text{words with } e \text{ ones}\}$
  with $\binom{k}{e}$ goes through indicator vectors of
  \texttt{Finset.powersetCard}. Axioms within \texttt{[propext,
  Classical.choice, Quot.sound]}; no \texttt{sorry}.
\item \texttt{CountingCorollary.lean} (same Lake project): \emph{the
  combinatorial core of Corollary~\ref{cor:count}}. With $T$ the
  compressed map and \texttt{stays m n} the predicate ``the first $m$
  iterates stay $\ge n$'' (which contains $\mathcal{E}\cap(2^m,2^{m+1}]$),
  \texttt{growth\_ineq} proves the integer form
  $2^{m(e+1)}\le(3\cdot 2^m+1)^e$ of $e\ge m\theta_M$ for $n>2^m$,
  \texttt{parityset\_injective} proves the Terras injectivity of
  $n\mapsto\{i<m: T^i n \text{ odd}\}$ on $(2^m,2^{m+1}]$ via the affine
  difference identity $2^j(T^jn-T^jn')=3^{e_j}(n-n')$, and
  \texttt{count\_le} concludes
  $\#\{n\in(2^m,2^{m+1}]:\texttt{stays}\}\le\sum_{e}\binom{m}{e}$ over the
  admissible $e$. Only the entropy estimate
  $\sum_{e\ge pm}\binom me\le 2^{mH(p)}$ and the passage from the integer
  inequality to $\theta_M$ are left to the text. Axioms within
  \texttt{[propext, Classical.choice, Quot.sound]}; no \texttt{sorry}.
\end{itemize}
Of the main results, Theorem~A (modulo the trivial reading of an
injection as a statement about distinct classes), the binomial
sandwich of Theorem~B, and the combinatorial core of
Corollary~\ref{cor:count} are thus machine-verified, while the entropy limit
in Theorem~B (Fekete and Stirling) and the remaining steps of A$'$, as
well as Theorems~C--G, are verified by conventional proof plus the
independent computations described above. The three core files use no
library at all and can be checked in seconds; the Mathlib project requires
a Mathlib cache. Both toolchains used (\texttt{v4.32.2}, released
2026-07-28, and \texttt{v4.33.1}) postdate the fix for the Lean kernel
soundness issue \#14576 of July 2026, which in any case was reachable only
through metaprogrammed kernel submissions; the artifacts use no
metaprogramming.

\subsection{Reproducibility and data availability}
All computations use exact integer (or exact rational) arithmetic. The
supplementary bundle contains the certificate artifact, the Lean file with
its pinned toolchain, the independent verifier, and one self-contained
script per computational claim of this paper, each printing the quantity it
certifies; a README lists the exact commands and expected outputs. SHA-256
hashes of the frozen artifacts:
\begin{center}
\small
\begin{tabular}{ll}
\toprule
artifact & SHA-256\\
\midrule
\texttt{certificates.json} & \texttt{d32a24e4b34eb33a0aa8f01406354ae2}\\
 & \texttt{2ddd65a7329b3d3309197150ad3d78fc}\\
\texttt{Barrier.lean} & \texttt{c416508cda41b0df105bae91801d70c3}\\
 & \texttt{7a3e2f65c0251bae1c90683492f4ec7d}\\
\texttt{Collapse.lean} & \texttt{08191fc7b0114d8504a9c930ae911ad7}\\
 & \texttt{075f33d05a38b6e26612bad2cf50b2f3}\\
\texttt{Rotation.lean} & \texttt{20f91338f3d623a496ce1d0a6b1ff775}\\
 & \texttt{75f14921756e17dc13ab48bd5863b97a}\\
\texttt{Counting.lean} & \texttt{09619f3b7862934a14ef7fb9ed9d4646}\\
 & \texttt{a50ea425a4622f819acf27bded5a4396}\\
\texttt{TheoremB.lean} & \texttt{0b2f7f732e6e42dc45624531a3e5c05a}\\
\texttt{CountingCorollary.lean} & \texttt{eb05554941149f2930d3579a984106b3}\\
 & \texttt{8847c88912e3ed3321a716dd8fd3879a}\\
\bottomrule
\end{tabular}
\end{center}

\section{Discussion: what is and is not blocked}\label{sec:discussion}

Theorems~A and A$'$ block the finite-horizon descent-certificate strategy
in full generality --- any mixture of depths, congruence scopes of arbitrary
moduli, disjunctions, adaptive stopping, and finite exceptions --- and
Theorems~F--G block its bounded-memory Lyapunov and finite-window
relaxations. It is equally important to state what is \emph{not} blocked.

\emph{Congruence information is exhausted.} By Theorem~A$'$, no
repackaging of congruence conditions on the starting point --- dyadic,
$3$-adic, or mixed --- escapes the count: the certificate frontier is
canonical, and its cost $2^{H(\theta)k}$ is invariant. A proof of the
conjecture by descent certificates would need scopes that are not
eventually periodic sets: sets defined by inequalities on orbit values, or
by information of unbounded window length. The same collapse yields
Corollary~\ref{cor:count} --- unconditionally at most $N^{0.950}$
potential counterexamples up to $N$ --- and, since the congruence
necessary conditions are attained exactly, improving that exponent
requires non-shadow information.

\emph{Unbounded memory and Lyapunov potentials.} Theorem~F excludes only
per-step, bounded-memory dyadic corrections, and Theorem~A$'$ closes
variable stopping for the potential $\log_2 n$ itself. Lyapunov schemes
with unbounded memory, or with potentials other than $\log_2 n$ certified
at variable times, are untouched --- though we note that the certificate
tree itself is such an unbounded-memory object, and Theorem~A prices its frontier at
$2^{H(\theta)k}$ classes per depth $k$.

\emph{Non-dyadic information along the orbit.} Theorems~F--G are
$2$-adic, and Theorem~A$'$ blocks $3$-adic conditions only as congruence
scopes on the initial point. Methods that consume $3$-adic
equidistribution \emph{along the orbit of an ensemble} --- in particular
the Fourier-analytic route of Tao \cite{Tao2022}, which transports the
initial distribution of an ensemble rather than finite windows of a single
orbit --- remain outside the scope of all our obstructions. We analyze the
quantitative $3$-adic mixing problem in a companion paper.

\emph{The cycle route.} Diophantine bounds on nontrivial cycles (linear
forms in logarithms combined with verified lower bounds on cycle elements
\cite{Eliahou1993,SimonsDeWeger2005,Rhin1987,Barina2021}) address the
periodicity half of the conjecture and are independent of everything here;
the appendix records the standard bound in the form we use.

Finally, the barrier perspective suggests reading the classical density
theorems not as progress toward the conjecture but as exact complexity
measurements of a proof system: the certified density $1-2^{-(1-H(\theta))k}$
and the certificate count $2^{H(\theta)k}$ are two faces of the same
entropy, and the fine structure of the frontier is governed by the
continued fraction of $\log_2 3$ (Theorems~C--E) --- the same Diophantine
data that governs the cycle route. We take this coincidence as evidence
that the difficulty quantified here is intrinsic to the problem and not to
one formalization.

\section*{Funding}

No funding was obtained for this work.

\section*{Disclosure statement}

The author reports there are no competing interests to declare.

\section*{Declaration of generative AI use}

Generative AI (Anthropic's Claude) was used, under the author's direction,
to assist in drafting and revising the manuscript, in developing and
checking the computational scripts and the Lean~4 proof files, and in
searching the literature. All definitions, theorems, and proofs were
reviewed by the author, the computational claims are backed by the
verifiable artifacts described in this article, and the author takes full
responsibility for its content.

\section*{Data availability statement}
All code and data supporting this study --- the certificate artifact, the
Lean~4 proof with pinned toolchain, the independent verifier, and one
self-contained exact-arithmetic script per computational claim, with a
claim-to-script map and SHA-256 manifest --- are openly available in the
reproduction repository \texttt{collatz-barrier-repro}
(\url{https://github.com/WakauraH/collatz-barrier-repro}); a
DOI-stamped archival snapshot of the exact version used here will be
deposited upon acceptance. The same material is included as the
supplementary bundle of this submission.
 
 
\begin{thebibliography}{99}
   
\bibitem{Terras1976}
R.~Terras,
\emph{A stopping time problem on the positive integers},
Acta Arith. \textbf{30} (1976), 241--252.

\bibitem{Everett1977}
C.~J. Everett,
\emph{Iteration of the number-theoretic function $f(2n)=n$, $f(2n+1)=3n+2$},
Adv. Math. \textbf{25} (1977), 42--45.

\bibitem{Lagarias1985}
J.~C. Lagarias,
\emph{The $3x+1$ problem and its generalizations},
Amer. Math. Monthly \textbf{92} (1985), 3--23.

\bibitem{Lagarias2010}
J.~C. Lagarias (ed.),
\emph{The Ultimate Challenge: The $3x+1$ Problem},
American Mathematical Society, Providence, RI, 2010.

\bibitem{Tao2022}
T.~Tao,
\emph{Almost all orbits of the Collatz map attain almost bounded values},
Forum Math. Pi \textbf{10} (2022), e12.

\bibitem{Barina2021}
D.~Barina,
\emph{Convergence verification of the Collatz problem},
J. Supercomput. \textbf{77} (2021), 2681--2688.

\bibitem{Eliahou1993}
S.~Eliahou,
\emph{The $3x+1$ problem: new lower bounds on nontrivial cycle lengths},
Discrete Math. \textbf{118} (1993), 45--56.

\bibitem{SimonsDeWeger2005}
J.~Simons and B.~de~Weger,
\emph{Theoretical and computational bounds for $m$-cycles of the $3n+1$-problem},
Acta Arith. \textbf{117} (2005), 51--70.

\bibitem{Rhin1987}
G.~Rhin,
\emph{Approximants de Pad\'e et mesures effectives d'irrationalit\'e},
in: S\'eminaire de Th\'eorie des Nombres, Paris 1985--86,
Progr. Math. \textbf{71}, Birkh\"auser, 1987, 155--164.

\bibitem{KrasikovLagarias2003}
I.~Krasikov and J.~C. Lagarias,
\emph{Bounds for the $3x+1$ problem using difference inequalities},
Acta Arith. \textbf{109} (2003), 237--258.

\bibitem{ApplegateLagarias1995}
D.~Applegate and J.~C. Lagarias,
\emph{Density bounds for the $3x+1$ problem. I. Tree-search method},
Math. Comp. \textbf{64} (1995), 411--426.

\bibitem{Wirsching1998}
G.~J. Wirsching,
\emph{The Dynamical System Generated by the $3n+1$ Function},
Lecture Notes in Math. \textbf{1681}, Springer, Berlin, 1998.

\bibitem{Conway1972}
J.~H. Conway,
\emph{Unpredictable iterations},
in: Proc. 1972 Number Theory Conference, Univ. Colorado, Boulder, 1972,
49--52.

\bibitem{KurtzSimon2007}
S.~A. Kurtz and J.~Simon,
\emph{The undecidability of the generalized Collatz problem},
in: TAMC 2007, Lecture Notes in Comput. Sci. \textbf{4484}, Springer, 2007,
542--553.

\bibitem{BernsteinLagarias1996}
D.~J. Bernstein and J.~C. Lagarias,
\emph{The $3x+1$ conjugacy map},
Canad. J. Math. \textbf{48} (1996), 1154--1169.

\bibitem{Korec1994}
I.~Korec,
\emph{A density estimate for the $3x+1$ problem},
Math. Slovaca \textbf{44} (1994), 85--89.

\bibitem{RozierTerracol2026}
O.~Rozier and C.~Terracol,
\emph{Paradoxical behavior in Collatz sequences},
Discrete Math. (2026), art.~115167; arXiv:2502.00948.

\bibitem{Niu2026}
T.~Niu,
\emph{Parity vectors and paradoxical sequences in the accelerated Collatz map},
arXiv:2605.13886 (2026).

\bibitem{FernandezIbanez2026}
C.~Fern\'andez and S.~Ib\'a\~nez,
\emph{Christoffel words as extremal structures in Collatz dynamics},
arXiv:2607.24844 (2026).

\bibitem{Knight2026}
K.~Knight,
\emph{Collatz high cycles do not exist},
Discrete Math. \textbf{349} (2026), no.~3, art.~114812.

\bibitem{Sharpe2026}
M.~Sharpe,
\emph{The Collatz conjecture at the critical line: machine-verified no-go,
density, and cycle theorems},
preprint and Lean~4 development (June 2026),
\url{https://github.com/msharpe248/collatz}.

\bibitem{Lothaire2002}
M.~Lothaire,
\emph{Algebraic Combinatorics on Words},
Encyclopedia Math. Appl. \textbf{90}, Cambridge University Press, 2002.

\end{thebibliography}

\appendix

\section{The standard cycle bound (reproduction)}\label{app:cycle}

For completeness we record the cycle bound used in
Section~\ref{sec:discussion}; the method is that of Eliahou
\cite{Eliahou1993} with the current verification bound, and no novelty is
claimed. Let a nontrivial cycle of the odd-to-odd accelerated map have odd
elements $m_1,\dots,m_K$ with $2$-adic valuations summing to $A$.
Multiplying the relations $m_{j+1}2^{a_j}=3m_j+1$ around the cycle gives
$2^A=\prod_j(3+1/m_j)$, hence with $M=\min_j m_j$,
\[
0<A-K\log_2 3\le \frac{K}{3M\ln 2}.
\]
With $M>2^{68}$ \cite{Barina2021} the right side is below
$K\cdot 2^{-69}$, forcing $A=\ceil{K\log_2 3}$ and making
$\|K\log_2 3\|^{+}$ (the distance to the integer \emph{above}) at most
$K\cdot 2^{-69}$. Running the one-sided ladder of best approximations of
$\log_2 3$ (mediant steps between convergents, compared by exact integer
powers $2^p$ vs $3^q$) shows the smallest such $K$ satisfies
\[
K\ \ge\ 8{,}963{,}457{,}696,
\]
so any nontrivial cycle has more than $2.3\times 10^{10}$ steps of the map
$C$. Stronger bounds via $m$-cycle theory are known
\cite{SimonsDeWeger2005}, and high cycles have recently been eliminated
altogether \cite{Knight2026}.


\end{document}
